\documentclass[11pt,letterpaper]{article}
\usepackage{cornell-notes}
\usepackage{needspace}

\newcommand{\CornellDocumentTitle}{Chapter 85: energy infrastructure cyber security}
\title{\CornellDocumentTitle}
\author{}
\date{}
\setDocTitle{\CornellDocumentTitle}
\setDocSubtitle{Cybersecurity Cornell Notes}
\setCornellCollection{Cybersecurity Reference Notes}
\setCornellUnitType{Chapter}
\setCornellUnitNumber{85}
\setCornellUnitTitle{energy infrastructure cyber security}
\setCornellCompanionLabel{Cybersecurity study companion}

\begin{document}
\maketitle
\begin{CornellOverviewBox}{Chapter Orientation}
\textbf{Chapter orientation.} This chapter examines energy infrastructure cyber security within the broader domain of critical infrastructure security. Its central problem is how to reason about energy, act, national, and cyber while keeping security objectives, operating assumptions, evidence, and recovery connected. A practitioner should approach the material through physical consequences, safety, availability, environmental dependencies, remote connectivity, maintenance, and recovery. Useful prerequisites include basic risk, threat, control, and systems concepts.
\end{CornellOverviewBox}
\section*{Learning Objectives}
\begin{itemize}
\item Explain the security purpose and scope of Energy Infrastructure Cyber Security.
\item Identify the assets, actors, assumptions, and trust boundaries associated with energy, act, and national.
\item Distinguish Introduction from Energy Sector Cyber Security Program As It Relates To The Electric Power Grid And Pipelines Subsectors and explain where their responsibilities intersect.
\item Analyze how failures in energy, act, and national can affect confidentiality, integrity, availability, privacy, or safety.
\item Evaluate relevant controls through physical consequences, safety, availability, environmental dependencies, remote connectivity, maintenance, and recovery.
\item Audit an implementation using asset and dependency maps, physical-access records, sensor telemetry, safety constraints, maintenance logs, and recovery exercises.
\item Prioritize remediation according to exploitability, consequence, control strength, and recovery readiness.
\item Verify that corrective actions change observable risk rather than documentation alone.
\end{itemize}
\section*{Cornell Cue-and-Notes}
\Needspace{8\baselineskip}
\subsection*{Introduction}
\begin{longtable}{>{\RaggedRight\arraybackslash}p{0.29\textwidth}|>{\RaggedRight\arraybackslash}p{0.65\textwidth}}
\CornellHeader
\CornellNoteRow{What security problem does Introduction address?}{\textbf{Core idea.} Treat Introduction as a structured part of Energy Infrastructure Cyber Security, with particular attention to terest national, securing energy, infrastructure terest, and energy critical. \textbf{Analytical lens.} This topic establishes a component of the chapter's argument; connect its definitions and mechanisms to a testable security objective. For terest national, securing energy, and infrastructure terest, model safety and physical consequences alongside conventional information-security effects. \textbf{Security reasoning.} Identify what must remain protected, who or what can alter the expected state, which assumptions cross a trust boundary, and how failure changes confidentiality, integrity, availability, privacy, or safety. \textbf{Application.} The practitioner should evaluate cyber and physical failure together, enforce safe states, and test operational recovery. \textbf{Evidence.} Seek asset and dependency maps, physical-access records, sensor telemetry, safety constraints, maintenance logs, and recovery exercises.}
\end{longtable}
\begin{CornellSummaryBox}{Topic Summary}
Introduction is best remembered as the connection between terest national, securing energy, infrastructure terest, and energy critical and the chapter's larger control objective. A defensible review states the assumption, expected behavior, evidence, failure consequence, and required response.
\end{CornellSummaryBox}
\Needspace{8\baselineskip}
\subsection*{Energy Sector Cyber Security Program As It Relates To The Electric Power Grid And Pipelines Subsectors}
\begin{longtable}{>{\RaggedRight\arraybackslash}p{0.29\textwidth}|>{\RaggedRight\arraybackslash}p{0.65\textwidth}}
\CornellHeader
\CornellNoteRow{Which assumptions matter in Energy Sector Cyber Security Program As It Relates To The Electric Power Grid And Pipelines Subsectors?}{\textbf{Core idea.} Treat Energy Sector Cyber Security Program As It Relates To The Electric Power Grid And Pipelines Subsectors as a structured part of Energy Infrastructure Cyber Security, with particular attention to energy, cyber, control, and critical. \textbf{Analytical lens.} This topic establishes a component of the chapter's argument; connect its definitions and mechanisms to a testable security objective. For energy, cyber, and control, model safety and physical consequences alongside conventional information-security effects. \textbf{Security reasoning.} Identify what must remain protected, who or what can alter the expected state, which assumptions cross a trust boundary, and how failure changes confidentiality, integrity, availability, privacy, or safety. \textbf{Application.} The practitioner should evaluate cyber and physical failure together, enforce safe states, and test operational recovery. \textbf{Evidence.} Seek asset and dependency maps, physical-access records, sensor telemetry, safety constraints, maintenance logs, and recovery exercises. Related subtopics include Industrial Control System Security Risks.}
\end{longtable}
\begin{CornellSummaryBox}{Topic Summary}
Energy Sector Cyber Security Program As It Relates To The Electric Power Grid And Pipelines Subsectors is best remembered as the connection between energy, cyber, control, and critical and the chapter's larger control objective. A defensible review states the assumption, expected behavior, evidence, failure consequence, and required response.
\end{CornellSummaryBox}
\Needspace{8\baselineskip}
\subsection*{Summary Of Risks To Computer-Based Control Systems Used In Energy Sector Subsectors And Recent Cyber Attacks}
\begin{longtable}{>{\RaggedRight\arraybackslash}p{0.29\textwidth}|>{\RaggedRight\arraybackslash}p{0.65\textwidth}}
\CornellHeader
\CornellNoteRow{How should a reviewer evaluate Summary Of Risks To Computer-Based Control Systems Used In Energy Sector Subsectors And Recent Cyber Attacks?}{\textbf{Core idea.} Treat Summary Of Risks To Computer-Based Control Systems Used In Energy Sector Subsectors And Recent Cyber Attacks as a structured part of Energy Infrastructure Cyber Security, with particular attention to grid, control, power, and scada. \textbf{Analytical lens.} This topic is threat-centered: separate actor capability, reachable attack surface, prerequisites, execution path, and consequence. For grid, control, and power, model safety and physical consequences alongside conventional information-security effects. \textbf{Security reasoning.} Identify what must remain protected, who or what can alter the expected state, which assumptions cross a trust boundary, and how failure changes confidentiality, integrity, availability, privacy, or safety. \textbf{Application.} The practitioner should evaluate cyber and physical failure together, enforce safe states, and test operational recovery. \textbf{Evidence.} Seek asset and dependency maps, physical-access records, sensor telemetry, safety constraints, maintenance logs, and recovery exercises. Related subtopics include Cyber Security Risks.}
\end{longtable}
\begin{CornellSummaryBox}{Topic Summary}
Summary Of Risks To Computer-Based Control Systems Used In Energy Sector Subsectors And Recent Cyber Attacks is best remembered as the connection between grid, control, power, and scada and the chapter's larger control objective. A defensible review states the assumption, expected behavior, evidence, failure consequence, and required response.
\end{CornellSummaryBox}
\Needspace{8\baselineskip}
\subsection*{Review Of Legislative Authorities And Policy Guidance Directing Us Department Of Energy Cybersecurity Programs And Activities And It'S Relationship With National}
\begin{longtable}{>{\RaggedRight\arraybackslash}p{0.29\textwidth}|>{\RaggedRight\arraybackslash}p{0.65\textwidth}}
\CornellHeader
\CornellNoteRow{Why can Review Of Legislative Authorities And Policy Guidance Directing Us Department Of Energy Cybersecurity Programs And Activities And It'S Relationship With National fail in practice?}{\textbf{Core idea.} Treat Review Of Legislative Authorities And Policy Guidance Directing Us Department Of Energy Cybersecurity Programs And Activities And It'S Relationship With National as a structured part of Energy Infrastructure Cyber Security, with particular attention to act, energy, national, and infrastructure. \textbf{Analytical lens.} This topic is governance-centered: connect required intent to accountable ownership, implementation, evidence, exceptions, and corrective action. For act, energy, and national, model safety and physical consequences alongside conventional information-security effects. \textbf{Security reasoning.} Identify what must remain protected, who or what can alter the expected state, which assumptions cross a trust boundary, and how failure changes confidentiality, integrity, availability, privacy, or safety. \textbf{Application.} The practitioner should evaluate cyber and physical failure together, enforce safe states, and test operational recovery. \textbf{Evidence.} Seek asset and dependency maps, physical-access records, sensor telemetry, safety constraints, maintenance logs, and recovery exercises. Related subtopics include Laboratories And Other, Federal Agencies With, Cybersecurity Roles, and Summary of Legislative Authorities,.}
\end{longtable}
\begin{CornellSummaryBox}{Topic Summary}
Review Of Legislative Authorities And Policy Guidance Directing Us Department Of Energy Cybersecurity Programs And Activities And It'S Relationship With National is best remembered as the connection between act, energy, national, and infrastructure and the chapter's larger control objective. A defensible review states the assumption, expected behavior, evidence, failure consequence, and required response.
\end{CornellSummaryBox}
\Needspace{8\baselineskip}
\subsection*{Concluding Discussion Of Key Policy Issues For Congress}
\begin{longtable}{>{\RaggedRight\arraybackslash}p{0.29\textwidth}|>{\RaggedRight\arraybackslash}p{0.65\textwidth}}
\CornellHeader
\CornellNoteRow{What security problem does Concluding Discussion Of Key Policy Issues For Congress address?}{\textbf{Core idea.} Treat Concluding Discussion Of Key Policy Issues For Congress as a structured part of Energy Infrastructure Cyber Security, with particular attention to internet, infrastructure, grid, and critical. \textbf{Analytical lens.} This topic is cryptography-centered: state the intended property and validate algorithms, parameters, keys, protocol binding, and lifecycle controls. For internet, infrastructure, and grid, model safety and physical consequences alongside conventional information-security effects. \textbf{Security reasoning.} Identify what must remain protected, who or what can alter the expected state, which assumptions cross a trust boundary, and how failure changes confidentiality, integrity, availability, privacy, or safety. \textbf{Application.} The practitioner should evaluate cyber and physical failure together, enforce safe states, and test operational recovery. \textbf{Evidence.} Seek asset and dependency maps, physical-access records, sensor telemetry, safety constraints, maintenance logs, and recovery exercises.}
\end{longtable}
\begin{CornellSummaryBox}{Topic Summary}
Concluding Discussion Of Key Policy Issues For Congress is best remembered as the connection between internet, infrastructure, grid, and critical and the chapter's larger control objective. A defensible review states the assumption, expected behavior, evidence, failure consequence, and required response.
\end{CornellSummaryBox}
\begin{CornellWarningBox}{Historical Context}
Some products, protocols, examples, threat observations, or legal references in this chapter are historically bounded. Retain the underlying threat model and control objective, but verify present applicability before treating a named implementation as current guidance.
\end{CornellWarningBox}
\begin{CornellOverviewBox}{Defensive Guidance}
Apply the chapter by making assets, authority, trust, expected behavior, and failure response explicit. In practice, evaluate cyber and physical failure together, enforce safe states, and test operational recovery. A control is not fully supported until its operation, monitoring, ownership, and recovery behavior are demonstrated with evidence.
\end{CornellOverviewBox}
\section*{Threat-and-Control Map}
\begingroup\scriptsize\setlength{\tabcolsep}{2pt}
\begin{longtable}{>{\RaggedRight\arraybackslash}p{0.135\textwidth}>{\RaggedRight\arraybackslash}p{0.145\textwidth}>{\RaggedRight\arraybackslash}p{0.155\textwidth}>{\RaggedRight\arraybackslash}p{0.145\textwidth}>{\RaggedRight\arraybackslash}p{0.16\textwidth}>{\RaggedRight\arraybackslash}p{0.17\textwidth}}
\toprule\rowcolor{Navy}\color{white}\textbf{Asset} & \color{white}\textbf{Threat / Failure} & \color{white}\textbf{Vulnerability} & \color{white}\textbf{Impact} & \color{white}\textbf{Detection / Evidence} & \color{white}\textbf{Control}\\\midrule
\endfirsthead
\toprule\rowcolor{Navy}\color{white}\textbf{Asset} & \color{white}\textbf{Threat / Failure} & \color{white}\textbf{Vulnerability} & \color{white}\textbf{Impact} & \color{white}\textbf{Detection / Evidence} & \color{white}\textbf{Control}\\\midrule
\endhead
People, physical processes, connected devices, and essential services & Tampering, unsafe command, service disruption, surveillance, or cascading failure & Insecure remote access, fragile dependencies, weak update paths, or insufficient safety separation & Safety events, prolonged outage, physical damage, privacy harm, or public disruption & Operational telemetry, integrity monitoring, physical controls, and anomaly investigation & Layered cyber-physical protection, safe-state design, segmentation, secure maintenance, and rehearsed recovery\\\midrule
Assumptions surrounding energy & Misuse or unexpected state transition & Trust in act is not validated & A local weakness becomes a broader compromise path & Review evidence for national and test negative paths & Make trust explicit, constrain authority, and fail safely\\\midrule
Recovery capability and decision evidence & Control failure remains unnoticed or cannot be reversed & Monitoring, ownership, or restoration has not been exercised & Longer exposure and slower, less reliable recovery & Alert-to-response exercises and restoration tests & Assigned ownership, actionable telemetry, preserved evidence, and rehearsed recovery\\\midrule
\bottomrule
\end{longtable}
\endgroup
\section*{Practical Security Checklist}
\begin{itemize}[label=$\square$]
\item Define the in-scope assets, users, services, data, and operational dependencies for Energy Infrastructure Cyber Security.
\item Document the expected behavior and security objective of energy.
\item Map identities, privileges, trust boundaries, and externally controlled inputs associated with energy and act.
\item Record the threat or failure being addressed, its prerequisites, likely path, and credible consequences.
\item Inspect asset and dependency maps, physical-access records, sensor telemetry, safety constraints, maintenance logs, and recovery exercises.
\item Test both the intended path and negative paths, including invalid state, partial failure, excessive privilege, and unavailable dependencies.
\item Confirm preventive controls are complemented by actionable detection and a named response owner.
\item Exercise containment, restoration, and national; record actual recovery behavior and unresolved dependencies.
\item Validate exceptions, compensating controls, expiration dates, and accountable approval.
\item Preserve reproducible evidence and tie each finding to affected assets, risk, remediation, and verification criteria.
\end{itemize}
\begin{CornellSummaryBox}{Chapter Synthesis}
The chapter's principal lesson is that energy infrastructure cyber security must be evaluated as an operating security system rather than an isolated product or statement. The most important failure patterns involve energy, act, national, and cyber, especially when ownership, boundaries, telemetry, or recovery remain implicit. A reviewer should connect architecture and behavior to asset and dependency maps, physical-access records, sensor telemetry, safety constraints, maintenance logs, and recovery exercises, prioritize findings by reachable consequence and control independence, and verify remediation through observable change. Within Critical Infrastructure Security, this chapter contributes a reusable method: state the objective, expose assumptions, test failure, preserve evidence, and learn from operation.
\end{CornellSummaryBox}
\section*{Self-Test Questions}
\begin{enumerate}
\item What is the central security objective of Energy Infrastructure Cyber Security?
\item How does Introduction contribute to the chapter's broader security argument?
\item Distinguish Energy Sector Cyber Security Program As It Relates To The Electric Power Grid And Pipelines Subsectors from Summary Of Risks To Computer-Based Control Systems Used In Energy Sector Subsectors And Recent Cyber Attacks.
\item Which trust assumptions should be made explicit when evaluating energy?
\item How could a weakness involving act affect confidentiality, integrity, availability, privacy, or safety?
\item What evidence would demonstrate that controls surrounding national operate as intended?
\item Why should prevention, detection, response, and recovery be evaluated together in this domain?
\item Scenario: A connected physical process remains functionally available after a cyber event but no longer operates within safe constraints. What should the reviewer examine first, and why?
\item Scenario: A control associated with cyber passes a configuration check but fails during an operational exercise. How should the finding be interpreted and prioritized?
\item How should a practitioner distinguish enduring principles in this chapter from time-sensitive tools, examples, or implementation details?
\end{enumerate}
\Needspace{8\baselineskip}
\section*{Answer Key}
\begin{enumerate}
\item The objective is to protect the relevant assets by understanding physical consequences, safety, availability, environmental dependencies, remote connectivity, maintenance, and recovery and by connecting controls to measurable risk reduction.
\item Introduction provides one part of the causal chain: it should identify expected behavior, dependencies, failure modes, and the evidence needed to justify confidence.
\item Energy Sector Cyber Security Program As It Relates To The Electric Power Grid And Pipelines Subsectors and Summary Of Risks To Computer-Based Control Systems Used In Energy Sector Subsectors And Recent Cyber Attacks address different portions of the problem. A strong comparison states their scope, assumptions, enforcement points, evidence, and how a failure in one affects the other.
\item Reviewers should identify who controls energy, what inputs or dependencies it trusts, where privilege changes, what happens during partial failure, and how those assumptions are tested.
\item A weakness in act can propagate across trust boundaries. The actual impact depends on reachable assets, granted authority, control independence, visibility, and recovery capability.
\item Useful evidence includes asset and dependency maps, physical-access records, sensor telemetry, safety constraints, maintenance logs, and recovery exercises; configuration alone is insufficient when operation, monitoring, or recovery is part of the claim.
\item Prevention reduces probability, detection limits dwell time, response limits spread, and recovery restores objectives. Omitting one layer creates an unexamined failure mode.
\item Begin by establishing scope, ownership, expected behavior, and the violated assumption. Then evaluate cyber and physical failure together, enforce safe states, and test operational recovery, preserving evidence for prioritization and follow-up.
\item Treat the exercise as stronger evidence than a static check. Prioritize according to reachable impact, likelihood of real operating conditions, missing compensating controls, and recovery difficulty.
\item Retain the principle, threat model, and control objective; separately label technologies, legal references, product examples, and threat observations whose applicability depends on time or environment.
\end{enumerate}
\section*{Key-Term Recap}
\begin{longtable}{>{\RaggedRight\arraybackslash}p{0.27\textwidth}>{\RaggedRight\arraybackslash}p{0.66\textwidth}}
\toprule\rowcolor{Navy}\color{white}\textbf{Term} & \color{white}\textbf{Meaning}\\\midrule
\endfirsthead
\toprule\rowcolor{Navy}\color{white}\textbf{Term} & \color{white}\textbf{Meaning}\\\midrule
\endhead
\textbf{Policy} & A management statement of required intent, responsibilities, and acceptable behavior. \\\midrule
\textbf{Threat} & A circumstance or actor capable of causing harm by exploiting a weakness or adverse condition. \\\midrule
\textbf{Risk} & The effect of uncertainty on objectives, commonly evaluated through likelihood, impact, exposure, and context. \\\midrule
\textbf{Resilience} & The ability to anticipate, withstand, recover from, and adapt to adverse conditions. \\\midrule
\textbf{Asset} & Information, service, capability, or physical resource whose value justifies protection. \\\midrule
\textbf{Authorization} & The decision process that determines which authenticated subject may perform which action on a resource. \\\midrule
\textbf{Incident Response} & The coordinated preparation, detection, analysis, containment, eradication, recovery, and learning activities for security incidents. \\\midrule
\textbf{Vulnerability} & A weakness that can be exploited or triggered to violate a security objective. \\\midrule
\bottomrule
\end{longtable}
\begin{CornellExamBox}{One-Sentence Takeaway}
Treat energy infrastructure cyber security as a verifiable chain from assets and assumptions through controls, evidence, response, and recovery.
\end{CornellExamBox}
\end{document}
